\chapter{Математическая модель решения задачи}\label{ch:6}

\section{Архитектура решения в нотации C4}

Перед началом разработки приложения рассмотрим архитектуру будущего решения в виде нотации C4 - она представлена на рисунках 12-15. 

На рисунке 12 изображен уровень системного контекста - здесь отражено взаимодействие клиента с системой. Посредником в этом процессе является мессенджер, через интерфейс которого пользователю будет удобно взаимодействовать с системой.

\begin{figure}[h]
\centering
\includegraphics[scale=0.5]{c4_system_context}
\caption{Уровень системного контекста в нотации C4}
\end{figure}

\begin{figure}[h]
\includegraphics[scale=0.37]{c4_containers}
\caption{Уровень контейнеров в нотации C4}
\end{figure}

На рисунке 13 представлен контейнерный уровень. Здесь раскрываются подробности внутренней архитектуры системы оценки стоимости ремонтных работ. На схеме показана информация о следующих элементах системы и их взаимодействии:

\begin{itemize}
    \item Bot Service - принимает сообщения от пользователя мессенджера, преобразует их, отправляет системе и возвращает ответы пользователю
    \item Backend API - главный серверный компонент, в котором сосредоточена основная бизнес-логика:
        \begin{itemize}
            \item Управление сценарием обработки обращений и расчёт стоимости работ
            \item Взаимодействие с базами данных
            \item Взаимодействие с брокером сообщений
        \end{itemize}
    \item Nginx - обратный прокси-сервер, принимающий входящие запросы. Балансировщик нагрузки
    \item Message Broker (Kafka) - брокер сообщений, обеспечивающий отказоустойчивое взаимодействие между бэкендом и ML
    \item ML Worker Pool - пул обработчиков для анализа изображений
    \item PostgreSQL - реляционная база данных для хранения служебных данных - информации о сценариях, статусах, результатах и др.
    \item MinIO S3 - объектное хранилище для хранения исходных и аннотированных изображений
\end{itemize}

\begin{figure}[h]
\centering
\includegraphics[scale=0.4]{c4_backend_api}
\caption{Компонентный уровень C4 - контейнер Backend API}
\end{figure}

Рассмотрим более подробно изнутри архитектуру самых важных частей системы - контейнеров Backend API и ML Worker'а.

Переходя к компонентному уровню нотации C4, начнём с контейнера бэкенда, в котором сосредоточена большая часть бизнес-логики (рисунок 14). От Bot Service'а к нему поступают запросы, которые обрабатываются компонентом API, использующем библиотеку FastAPI. Внутри реализуется стейт-машина и управление полным сценарием от отправки запроса в сервис инференса до получения ответа, взаимодействия с базами данных. Внутри контейнера имеются компоненты, позволяющие отказоустойчиво обработать каждый сценарий и отправить ответ пользователю, среди них:
\begin{itemize}
    \item Outbox Flusher - компонент для отказоустойчивой отправки накопленных событий из таблицы outbox в брокер сообщений
    \item Heartbeat Watchdog - компонент мониторинга зависших сценариев
\end{itemize}

\begin{figure}[h]
\centering
\includegraphics[scale=0.5]{c4_ml_worker}
\caption{Компонентный уровень C4 - контейнер ML Worker}
\end{figure}

Бэкенд взаимодействует с ML Worker'ами через брокер сообщений. Конечный вердикт по стоимости и длительности работ выполняется при помощи компонента расчёта стоимости.

Детальный обзор ML контейнера представлен на рисунке 15. Он получает сообщение от брокера сообщений о появившихся сценариях и производит цикл действий по обработке фотографии вплоть до формирования финального аннотированного изображения, которое сохраняет в S3. ML Worker состоит из следующих компонентов:
\begin{itemize}
    \item Предобработчик изображений - отвечает за подготовку изображения к анализу
    \item Детектор деталей автомобиля на изображении - определяет области расположения деталей
    \item Модуль обрезки изображений - отвечает за обрезку отдельных фрагментов изображения (деталей) по найденным областям
    \item Детектор повреждений автомобиля на изображении - определяет области расположения повреждений на обрезанных областях деталей
    \item Постобработчик изображения - формирует итоговое размеченное изображение 
\end{itemize}

\section{Схема базы данных}

\begin{figure}[h]
\centering
\includegraphics[scale=0.6]{db_schema}
\caption{Схема базы данных PostgreSQL}
\end{figure}

В системе имеется два хранилища данных - S3 MinIO для изображений и PostgreSQL для хранения данных о каждом сценарии и изображении. Рассмотрим схему базы данных PostgreSQL подробнее. Она создаётся скриптом инициализации при первом запуске контейнера и содержит несколько таблиц:

\begin{itemize}
    \item repair\_requests - хранение информации о заявке на оценку стоимости: информация о чате с пользователем, статусе, режиме обработки, изображении.
    \item detected\_parts - информация о результате сегментации деталей - тип детали, уверенность, ограничивающий прямоугольник
    \item detected\_damages - информация о найденных повреждениях - тип повреждения, уверенность, маска на изображении
    \item pricing\_rules - таблица для этапа расчёта стоимости и длительности работ на основе параметров <деталь - повреждение>
    \item outbox\_events - доменные события для Kafka
\end{itemize}

Визуализированная схема базы данных PostgreSQL с показанными связями и типами данных представлена на рисунке 16.



\section{Архитектура и математическая модель ML-составляющей}

Рассмотрим более подробно архитектуру контейнера ML Worker. Процесс внутри представлен в виде пайплайна подзадач - общий процесс детекции и сегментации повреждений разделён на несколько этапов, главными из которых являются компоненты, требующие обучения моделей детекции и сегментации деталей и повреждений (Причины выбора такого подхода описаны в главе 1).

Для решения задачи анализа изображения целесообразно выбрать модель семейства YOLO. Полноценное сравнение и анализ архитектур моделей производились в главе 1, и текущий выбор обусловлен тем, что YOLOv8 представляет собой компромисс между скоростью обучения и внедрения, скоростью инференса, вычислительной эффективностью и качеством предсказания.

В обобщённом виде архитектуру YOLOv8-seg можно представить как совокупность трёх частей:
\begin{itemize}
    \item backbone - извлечение признаков из входного изображения
    \item neck - объединение признаков разных масштабов
    \item segment head - предсказание классов, координат объектов и коэффициентов масок
\end{itemize}

Для обучения модели используются изображения повреждённых автомобилей. Подготовка данных и обучение поддерживают механизмы аугментации, включая Mosaic - преобразование, при котором несколько изображений объединяются в одно для улучшения обобщающей способности модели.

\subsection{Модель детекции и сегментации деталей автомобиля}

На вход системы поступает изображение автомобиля:
\[
X \in \mathbb{R}^{H \times W \times 3},
\]
где $H$ --- высота изображения, $W$ --- ширина изображения, а третье измерение соответствует цветовым каналам.

Перед передачей изображения в модель выполняется предобработка: изображение приводится к размеру входа модели 640 х 640 пикселей, а значения пикселей масштабируются к диапазону [0,1] путём деления на 255. Согласно документации Ultralytics \cite{src22}, такая схема соответствует стандартному способу подготовки данных для YOLO-моделей.

Первый этап пайплайна - изображение автомобиля обрабатывается моделью детекции и сегментации деталей. Обученная модель сегментации отображает входное изображение в множество предсказаний, каждое из которых содержит bounding box, маску, класс и confidence.

Результатом работы модели деталей является множество:
\[
S_{\text{parts}} = \left\{(b_i, m_i, p_i, s_i)\right\}_{i=1}^{N},
\]
где:
\begin{itemize}
    \item $b_i$ - bounding box (ограничивающий прямоугольник) детали в координатах исходного изображения
    \item $m_i$ - маска детали
    \item $p_i$ - класс детали автомобиля
    \item $s_i \in [0,1]$ - значение confidence (уверенности модели в предсказании)
    \item $N$ - количество найденных деталей
\end{itemize}

Для каждой обнаруженной на фотографии детали модель оценивает и выдаёт вероятность принадлежности к одному из классов: дверь, переднее или заднее крыло и другим - полный список деталей был описан в главе 4. 


После получения предсказаний выполняется фильтрация по порогу уверенности:
\[
S'_{\text{parts}} =
\left\{
(b_i, m_i, p_i, s_i) \in S_{\text{parts}} \;\middle|\;
s_i \geq \tau_{\text{parts}}
\right\},
\]
где $\tau_{\text{parts}}$ - порог уверенности модели в предсказании.

Для устранения дублирующихся и пересекающихся предсказаний применяется non-maximum suppression (NMS). Для двух bounding box-ов $b_i$ и $b_j$ коэффициент пересечения определяется как Intersection over Union (IoU) - мера перекрытия двух областей, которая определяется отношением площади их пересечения к площади объединения:
\[
IoU(b_i, b_j) = \frac{|b_i \cap b_j|}{|b_i \cup b_j|}.
\]

Если
\[
IoU(b_i, b_j) > \lambda,
\]
то предсказание с меньшим значением confidence отклоняется. NMS применяется раздельно для каждого класса объектов.

Обучение модели деталей выполняется с использованием комбинированной функции потерь, подробнее описанной в разделе 6.2.3.

\subsection{Формирование обрезанных областей деталей автомобиля}

После детекции, сегментации и фильтрации деталей выполняется построение кропов (обрезанных областей изображения). Для каждой детали из полученного списка выбирается ограничивающий прямоугольник $b_i$, по которому из исходного изображения вырезается соответствующая область.
На этом этапе формируется локальная область изображения в виде прямоугольника, содержащая только одну деталь автомобиля. Такой подход снижает влияние фона и повышает точность поиска повреждений, а также позволяет рассматривать каждую деталь как отдельный объект анализа.

\subsection{Модель детекции и сегментации повреждений}

Для каждой детали автомобиля выполняется отдельный этап поиска повреждений. На вход второй модели поступает не полное изображение автомобиля, а кроп детали. Индекс $i$ обозначает номер детали, для которой выполняется анализ: модель применяется независимо к каждому кропу.

Результат работы модели для $i$-й детали можно записать как:
\[
S_{\text{damages}, i} =
\left\{(b_{ij}, m_{ij}, d_{ij}, s_{ij})\right\}_{j=1}^{M_i},
\]
где:
\begin{itemize}
    \item $b_{ij}$ - bounding box (ограничивающий прямоугольник) повреждения в координатах кропа детали
    \item $m_{ij}$ - маска повреждения
    \item $d_{ij}$ - класс повреждения
    \item $s_{ij}$ - значение confidence (уверенности модели в предсказании)
    \item $M_i$ - количество найденных повреждений на $i$-й детали.
\end{itemize}

Общее множество повреждений по поданному изображению формируется как объединение результатов по всем кропам:
\[
S_{\text{damages}} = \bigcup_{i=1}^{N'} S_{\text{damages}, i},
\]
где $N'$ --- количество деталей, прошедших фильтрацию на первом этапе.

Для обучения модели повреждений используется тот же класс архитектуры YOLOv8-seg, однако датасет строится уже на изображениях кропов деталей, преимущества такого подхода были описаны выше.

Для повреждений также используется комбинированная функция потерь, включающая локализацию, классификацию, распределительную регрессию и сегментацию:
\[
\mathcal{L} =
\lambda_{\text{box}}\mathcal{L}_{\text{box}} +
\lambda_{\text{cls}}\mathcal{L}_{\text{cls}} +
\lambda_{\text{dfl}}\mathcal{L}_{\text{dfl}} +
\lambda_{\text{seg}}\mathcal{L}_{\text{seg}},
\]
где:
\begin{itemize}
    \item $\mathcal{L}_{\text{box}}$ - ошибка локализации bounding box;
    \item $\mathcal{L}_{\text{cls}}$ - ошибка классификации;
    \item $\mathcal{L}_{\text{dfl}}$ - Distribution Focal Loss;
    \item $\mathcal{L}_{\text{seg}}$ - ошибка сегментации масок;
    \item $\lambda_{\text{box}}, \lambda_{\text{cls}}, \lambda_{\text{dfl}}, \lambda_{\text{seg}}$ - веса соответствующих компонент.
\end{itemize}

Complete IoU (CIoU) расширяет IoU тремя дополнительными штрафными членами:
\[
CIoU(b_{\text{pred}}, b_{\text{gt}}) = IoU(b_{\text{pred}}, b_{\text{gt}})
- \frac{\rho^2(c_{\text{pred}}, c_{\text{gt}})}{d^2}
- \alpha v,
\]
где $\rho(c_{\text{pred}}, c_{\text{gt}})$ - евклидово расстояние между центрами предсказанного и эталонного прямоугольников, $d$ - диагональ минимального охватывающего прямоугольника, а $v$ - штраф за несоответствие соотношений сторон, взвешенный коэффициентом $\alpha$.

Компонент локализации можно записать как:
\[
\mathcal{L}_{\text{box}} = 1 - CIoU(b_{\text{pred}}, b_{\text{gt}}),
\]
где $b_{\text{pred}}$ и $b_{\text{gt}}$ - предсказанный и эталонный bounding box соответственно.


\subsection{Формирование финального аннотированного изображения}

После получения результатов на всех кропах необходимо выполнить перенос найденных масок обратно на исходное изображение. $m_1, m_2, \dots, m_k$ - маски всех обнаруженных повреждений, приведённые к размеру исходного изображения. Объединённая маска определяется как поэлементный максимум:
\[
M = \max(m_1, m_2, \ldots, m_k).
\]

Цветовое представление маски формируется путём назначения каждому типу повреждения собственного цвета. Тогда композитное изображение можно представить как:
\[
X^{\text{comp}} = (1 - \alpha M)\odot X + \alpha M \odot C,
\]
где $\alpha$ - коэффициент прозрачности, $C$ - цветовая карта маски, а $\odot$ - поэлементное умножение.


\section{Математическая модель расчёта стоимости и длительности ремонтных работ}

После завершения автоматического анализа изображения или ручного ввода повреждений требуется определить ориентировочную стоимость и длительность ремонтных работ.

Пусть для каждого найденного повреждения известны:
\begin{itemize}
    \item Тип детали ($p_i$)
    \item Тип повреждения ($d_i$)
\end{itemize}

Для каждого повреждения стоимость и длительность определяются как табличные отображения - каждая пара (тип детали, тип повреждения) однозначно соответствует фиксированным значениям стоимости и длительности работ.

Итоговая стоимость ремонта вычисляется как:
\[
C = \sum_{i \in A} c_i,
\]

а итоговая длительность ремонта:
\[
T = \sum_{i \in A} t_i.
\]
где А - множество найденных повреждений.

Таблицы 5 и 6, которые были описаны в главе 4, позволяют однозначно задать стоимость и длительность работ для каждой пары "деталь - тип повреждения". Это делает расчёт прозрачным и позволяет изменять ценники и правила автосервиса без изменения алгоритма системы.

